\documentclass[11pt,a4paper]{article}

\usepackage[margin=1in]{geometry}
\usepackage{amsmath,amssymb}
\usepackage{booktabs}
\usepackage{longtable}
\usepackage{enumitem}
\usepackage{fancyhdr}
\usepackage{titlesec}
\usepackage{hyperref}
\usepackage{xcolor}
\usepackage{listings}
\usepackage{fancyvrb}

\hypersetup{
    colorlinks=true,
    linkcolor=blue!70!black,
    urlcolor=blue!60!black,
}

\DeclareMathOperator*{\argmin}{arg\,min}

\lstset{
    basicstyle=\small\ttfamily,
    breaklines=true,
    frame=single,
    framesep=3pt,
    rulecolor=\color{black!30},
    backgroundcolor=\color{black!3},
    columns=flexible,
    keepspaces=true,
    showstringspaces=false,
    tabsize=4,
}

\lstdefinelanguage{Rust}{
    morekeywords={let,mut,fn,pub,struct,enum,impl,trait,use,mod,crate,self,super,
        vec,for,in,if,else,match,return,where,type,const,static,ref,
        Some,None,Ok,Err,Result,Option,Vec,Box,dyn,async,await,move},
    morestring=[b]",
    morecomment=[l]{//},
    morecomment=[s]{/*}{*/},
    sensitive=true,
}

\lstdefinelanguage{Python}{
    morekeywords={import,from,class,def,return,self,None,True,False,as,if,else,for,in,with},
    morestring=[b]",
    morestring=[b]',
    morecomment=[l]{\#},
    sensitive=true,
}

\pagestyle{fancy}
\fancyhf{}
\fancyhead[L]{\textit{scry-learn Architecture Comparison}}
\fancyhead[R]{\thepage}
\renewcommand{\headrulewidth}{0.4pt}

\titleformat{\section}{\Large\bfseries}{}{0em}{}[\titlerule]
\titleformat{\subsection}{\large\bfseries}{}{0em}{}

\title{\vspace{-1cm}\textbf{scry-learn Architecture}\\[0.3em]
\large Comparison with linfa, smartcore, and scikit-learn}
\author{scry-learn Technical Reference}
\date{}

\begin{document}
\maketitle
\thispagestyle{fancy}

\tableofcontents
\newpage

%=============================================================================
\section{Data Representation}
%=============================================================================

\subsection{Primary Data Containers}

\begin{center}
\begin{tabular}{llll}
\toprule
\textbf{Library} & \textbf{Core Type} & \textbf{Backing Storage} & \textbf{Layout} \\
\midrule
scry-learn & \texttt{Dataset} & \texttt{Vec<Vec<f64>{>}} + \texttt{DenseMatrix} & Column-major \\
linfa & \texttt{DatasetBase<R, T>} & \texttt{ndarray::Array2<F>} & Row-major \\
smartcore & \texttt{DenseMatrix<T>} & \texttt{Vec<T>} flat buffer & Row-major \\
scikit-learn & numpy \texttt{ndarray} & C-contiguous buffer & Row-major \\
\bottomrule
\end{tabular}
\end{center}

\subsubsection*{scry-learn}

\begin{lstlisting}[language=Rust]
// Column-major: features[feature_idx][sample_idx]
let data = Dataset::new(
    vec![
        vec![1.0, 2.0, 3.0],  // feature 0: 3 samples
        vec![4.0, 5.0, 6.0],  // feature 1: 3 samples
    ],
    vec![0.0, 1.0, 0.0],      // target
    vec!["x1".into(), "x2".into()],
    "y",
);
\end{lstlisting}

The \texttt{Dataset} struct bundles features, target, metadata (feature names, target name, class labels), and optional backing stores:

\begin{itemize}[itemsep=2pt]
\item \texttt{DenseMatrix}: contiguous column-major \texttt{Vec<f64>} (\texttt{data[col * n\_rows + row]}) for cache-friendly column scans
\item \texttt{row\_major\_cache}: lazily computed flat buffer for row-oriented algorithms (KNN, neural nets)
\item \texttt{CscMatrix} / \texttt{CsrMatrix}: custom compressed sparse column/row formats
\end{itemize}

\subsubsection*{linfa}

\begin{lstlisting}[language=Rust]
use linfa::prelude::*;
use ndarray::array;

let features = array![[1.0, 4.0], [2.0, 5.0], [3.0, 6.0]];
let targets = array![0.0, 1.0, 0.0];
let dataset = Dataset::new(features, targets);
\end{lstlisting}

\texttt{DatasetBase<R, T>} is generic over the record type \texttt{R} (typically \texttt{Array2<f64>}) and target type \texttt{T} (can be \texttt{f64}, \texttt{usize}, \texttt{bool}, or label types). Hard dependency on the \texttt{ndarray} crate.

\subsubsection*{smartcore}

\begin{lstlisting}[language=Rust]
use smartcore::linalg::basic::matrix::DenseMatrix;

let x = DenseMatrix::from_2d_array(&[
    &[1.0, 4.0], &[2.0, 5.0], &[3.0, 6.0]
]);
let y = vec![0.0, 1.0, 0.0];
\end{lstlisting}

\texttt{DenseMatrix<T>} is row-major (\texttt{data[row * n\_cols + col]}). Also supports \texttt{ndarray::Array2} and \texttt{nalgebra::DMatrix} via the \texttt{BaseMatrix} trait abstraction. No built-in feature names or metadata.

\subsubsection*{scikit-learn}

\begin{lstlisting}[language=Python]
import numpy as np
X = np.array([[1.0, 4.0], [2.0, 5.0], [3.0, 6.0]])
y = np.array([0.0, 1.0, 0.0])
\end{lstlisting}

No unified dataset struct. \texttt{X} and \texttt{y} are separate numpy arrays passed independently. \texttt{feature\_names\_in\_} attribute populated after \texttt{fit()} (since v1.0).

%-----------------------------------------------------------------------------
\subsection{Memory Layout Tradeoffs}

scry-learn is the only library in this comparison that defaults to \textbf{column-major} storage:

\textbf{Column-major advantages:}
\begin{itemize}[itemsep=2pt]
\item Tree split evaluation scans a single feature column---column-major keeps this contiguous in cache
\item Coordinate descent (Lasso, ElasticNet) processes one feature at a time
\item Variance/mean computation per feature is cache-friendly
\item CSC sparse format aligns naturally with column-major dense layout
\end{itemize}

\textbf{Column-major disadvantages:}
\begin{itemize}[itemsep=2pt]
\item Row-oriented algorithms (KNN, neural nets) require transposition, mitigated by the lazy \texttt{row\_major\_cache}
\item Prediction input is row-major (\texttt{\&[Vec<f64>]}), creating an API asymmetry
\end{itemize}

\textbf{Row-major advantages (linfa, smartcore, sklearn):}
\begin{itemize}[itemsep=2pt]
\item Each sample is contiguous---natural for per-sample operations
\item Consistent with numpy/BLAS conventions
\end{itemize}

Note: sklearn internally converts to Fortran-order (column-major) in many tree implementations for split finding, arriving at a similar layout to scry-learn's default.

%-----------------------------------------------------------------------------
\subsection{Sparse Data Support}

\begin{center}
\begin{tabular}{lll}
\toprule
\textbf{Library} & \textbf{Sparse Formats} & \textbf{Implementation} \\
\midrule
scry-learn & CSC, CSR & Custom (\texttt{sparse.rs}), zero external deps \\
linfa & None built-in & External \texttt{sprs} crate \\
smartcore & None & --- \\
scikit-learn & CSR, CSC, COO, BSR, LIL, DOK & \texttt{scipy.sparse} \\
\bottomrule
\end{tabular}
\end{center}

scry-learn implements CSC and CSR from scratch. \texttt{CscMatrix} is used for fitting (column-oriented coordinate descent, tree splits), \texttt{CsrMatrix} for prediction (row-oriented dot products). Algorithms auto-dispatch to sparse kernels when \texttt{data.sparse\_csc()} returns \texttt{Some}.

%-----------------------------------------------------------------------------
\subsection{Type System}

\begin{center}
\begin{tabular}{lll}
\toprule
\textbf{Library} & \textbf{Numeric Type} & \textbf{Flexibility} \\
\midrule
scry-learn & \texttt{f64} only & Simplest API, no generic parameters \\
linfa & \texttt{F: Float} (f32/f64) & Compile-time precision choice \\
smartcore & \texttt{T: RealNumber} & Generic over numeric types \\
scikit-learn & numpy dtype & Runtime dtype flexibility \\
\bottomrule
\end{tabular}
\end{center}

scry-learn's \texttt{f64}-only approach eliminates generic type parameters from all APIs, resulting in simpler function signatures and error messages.


%=============================================================================
\section{API Design Patterns}
%=============================================================================

\subsection{Model Lifecycle}

\subsubsection*{scry-learn: Mutable self, builder pattern}

\begin{lstlisting}[language=Rust]
let mut model = RandomForestClassifier::new()
    .n_estimators(100)
    .max_depth(10);

model.fit(&train)?;
let preds = model.predict(&test_features)?;
\end{lstlisting}

\begin{itemize}[itemsep=2pt]
\item \texttt{fit(\&mut self, data: \&Dataset) -> Result<()>} --- mutates in-place
\item \texttt{predict(\&self, features: \&[Vec<f64>]) -> Result<Vec<f64>{>}} --- immutable borrow
\item Runtime \texttt{NotFitted} error if predict called before fit
\end{itemize}

\subsubsection*{linfa: Type-state pattern, immutable models}

\begin{lstlisting}[language=Rust]
let model = DecisionTree::params()
    .max_depth(Some(10))
    .fit(&dataset)?;  // returns NEW fitted model

let predictions = model.predict(&test);
\end{lstlisting}

Unfitted params struct and fitted model are \textbf{different types} (e.g., \texttt{DecisionTreeParams} $\to$ \texttt{DecisionTree}). This provides a \textbf{compile-time guarantee} that predict cannot be called on an unfitted model.

\subsubsection*{smartcore: Static method, no unfitted state}

\begin{lstlisting}[language=Rust]
let model = RandomForestClassifier::fit(
    &x, &y, Default::default()
)?;
let predictions = model.predict(&x_test)?;
\end{lstlisting}

\subsubsection*{scikit-learn: Mutable self}

\begin{lstlisting}[language=Python]
model = RandomForestClassifier(n_estimators=100, max_depth=10)
model.fit(X_train, y_train)
predictions = model.predict(X_test)
\end{lstlisting}

%-----------------------------------------------------------------------------
\subsection{Comparison Summary}

\begin{center}
\begin{tabular}{lllll}
\toprule
\textbf{Aspect} & \textbf{scry-learn} & \textbf{linfa} & \textbf{smartcore} & \textbf{sklearn} \\
\midrule
Fit mutates self & Yes & No (new type) & No (static) & Yes \\
Compile-time fit check & No & Yes & Yes & No \\
Shared trait & \texttt{PipelineModel} & \texttt{Fit}/\texttt{Predict} & None & \texttt{BaseEstimator} \\
Config style & Builder methods & \texttt{with\_*} params & Params struct & Constructor kwargs \\
X/y bundling & \texttt{Dataset} & \texttt{DatasetBase} & Separate & Separate \\
\bottomrule
\end{tabular}
\end{center}

%-----------------------------------------------------------------------------
\subsection{Error Handling}

\begin{lstlisting}[language=Rust]
// scry-learn: typed error enum
#[non_exhaustive]
pub enum ScryLearnError {
    EmptyDataset,
    NotFitted,
    ShapeMismatch { expected: usize, got: usize },
    InvalidParameter(String),
    ConvergenceFailure(String),
    SchemaVersionMismatch { model: u32, current: u32 },
}
\end{lstlisting}

\begin{center}
\begin{tabular}{lll}
\toprule
\textbf{Library} & \textbf{Error Type} & \textbf{Pattern} \\
\midrule
scry-learn & \texttt{ScryLearnError} enum & \texttt{Result<T, ScryLearnError>} \\
linfa & Per-crate error types & \texttt{Result<T, linfa\_trees::Error>} \\
smartcore & \texttt{Failed} & \texttt{Result<T, Failed>} \\
scikit-learn & Python exceptions & \texttt{ValueError}, \texttt{NotFittedError} \\
\bottomrule
\end{tabular}
\end{center}

scry-learn uses a single error enum across all algorithms, simplifying error handling in pipelines. linfa's per-crate errors are more granular but require mapping when combining algorithms from different crates.

%-----------------------------------------------------------------------------
\subsection{Trait System}

\begin{center}
\begin{tabular}{lp{10cm}}
\toprule
\textbf{Library} & \textbf{Architecture} \\
\midrule
scry-learn & Concrete types. \texttt{PipelineModel} via \texttt{impl\_pipeline\_model!} macro. \texttt{Transformer} for preprocessors. \texttt{PartialFit} for online learning. \texttt{Tunable} for hyperparameter search. \\
linfa & Rich trait hierarchy: \texttt{Fit}, \texttt{FitWith}, \texttt{Predict}, \texttt{PredictInplace}, \texttt{Transformer}. \\
smartcore & Concrete types with inherent methods. No shared traits. \\
scikit-learn & Duck typing with \texttt{BaseEstimator} mixin. \texttt{RegressorMixin}/\texttt{ClassifierMixin} add \texttt{score()}. \\
\bottomrule
\end{tabular}
\end{center}


%=============================================================================
\section{Algorithm Implementation Comparison}
%=============================================================================

\subsection{Solver Choices}

{\small
\begin{center}
\begin{tabular}{lllll}
\toprule
\textbf{Algorithm} & \textbf{scry-learn} & \textbf{linfa} & \textbf{smartcore} & \textbf{sklearn} \\
\midrule
Linear Reg. & Normal, QR, SVD & SVD & SVD & SVD, Cholesky \\
Ridge & Normal + $\alpha I$ & SVD + $\alpha I$ & SVD & 6 solvers \\
Logistic Reg. & L-BFGS, GD & IRLS & LBFGS & 6 solvers \\
Lasso & Coord. descent & Coord. descent & Coord. descent & Cython CD \\
ElasticNet & Coord. descent & Coord. descent & --- & Cython CD \\
SVM (kernel) & Simplified SMO & SMO & SMO & libsvm \\
SVM (linear) & Pegasos SGD & --- & --- & liblinear \\
Decision Tree & CART, pre-sorted & CART & CART & Cython CART \\
Random Forest & Parallel (rayon) & Ensemble & Ensemble & joblib+Cython \\
Grad. Boosting & CART + Newton & --- & --- & CART + hist. \\
Hist-GBT & 256-bin, GPU & --- & --- & Cython hist. \\
K-Means & Lloyd's + k-means++ & Lloyd's & Lloyd's & Lloyd's, Elkan \\
PCA & Jacobi rotation & SVD & SVD & LAPACK SVD \\
KNN & Brute + KD-tree & Ball/KD-tree & Brute, cover & Brute, ball, KD \\
\bottomrule
\end{tabular}
\end{center}
}

%-----------------------------------------------------------------------------
\subsection{Parallelism}

\begin{center}
\begin{tabular}{lll}
\toprule
\textbf{Library} & \textbf{Framework} & \textbf{Scope} \\
\midrule
scry-learn & \texttt{rayon} & RF training, batch predict, CV, histograms \\
linfa & \texttt{rayon} (optional) & Select algorithms \\
smartcore & None & Single-threaded throughout \\
scikit-learn & \texttt{joblib} + OpenMP & Most ensembles, CV, KNN \\
\bottomrule
\end{tabular}
\end{center}

scry-learn's Random Forest uses rayon for parallel tree fitting with atomic OOB vote accumulation---each tree's OOB predictions are aggregated via \texttt{AtomicU64} arrays shared across threads.

%-----------------------------------------------------------------------------
\subsection{GPU Acceleration}

\begin{center}
\begin{tabular}{lll}
\toprule
\textbf{Library} & \textbf{GPU Support} & \textbf{Technology} \\
\midrule
scry-learn & Feature-gated (\texttt{gpu}) & wgpu compute shaders \\
linfa & None & --- \\
smartcore & None & --- \\
scikit-learn & None native & cuML (RAPIDS) separate \\
\bottomrule
\end{tabular}
\end{center}

scry-learn's GPU acceleration covers:
\begin{itemize}[itemsep=2pt]
\item \texttt{matmul}: Neural net forward/backward (threshold: \texttt{batch * max\_dim >= 4096})
\item \texttt{xtx\_xty}: Linear regression Gram matrix
\item \texttt{pairwise\_distances\_squared}: KNN distance computation
\item \texttt{build\_histograms}: Histogram GBT bin accumulation
\end{itemize}

The \texttt{ComputeBackend} trait abstracts CPU vs GPU, with \texttt{accel::auto()} selecting at runtime.


%=============================================================================
\section{Preprocessing Pipeline}
%=============================================================================

\subsection{Available Transformers}

{\small
\begin{center}
\begin{tabular}{lcccc}
\toprule
\textbf{Transformer} & \textbf{scry-learn} & \textbf{linfa} & \textbf{smartcore} & \textbf{sklearn} \\
\midrule
StandardScaler & Yes & Yes & Yes & Yes \\
MinMaxScaler & Yes & Yes & Yes & Yes \\
RobustScaler & Yes & --- & --- & Yes \\
Normalizer & Yes & --- & --- & Yes \\
PCA & Yes (Jacobi) & Yes (SVD) & Yes (SVD) & Yes \\
PolynomialFeatures & Yes & --- & --- & Yes \\
LabelEncoder & Yes & --- & --- & Yes \\
OneHotEncoder & Yes & --- & --- & Yes \\
SimpleImputer & Yes & --- & --- & Yes \\
ColumnTransformer & Yes & --- & --- & Yes \\
CountVectorizer & Yes & --- & --- & Yes \\
TfidfVectorizer & Yes & --- & --- & Yes \\
SelectKBest & Yes & --- & --- & Yes \\
VarianceThreshold & Yes & --- & --- & Yes \\
\bottomrule
\end{tabular}
\end{center}
}

%-----------------------------------------------------------------------------
\subsection{Pipeline Composition}

\begin{lstlisting}[language=Rust]
// scry-learn
let pipeline = Pipeline::new()
    .add_transformer(StandardScaler::new())
    .set_model(RandomForestClassifier::new());
pipeline.fit(&train)?;
let preds = pipeline.predict(&test)?;
\end{lstlisting}

\begin{center}
\begin{tabular}{llll}
\toprule
\textbf{Library} & \textbf{Pipeline} & \textbf{ColumnTransformer} & \textbf{FeatureUnion} \\
\midrule
scry-learn & \texttt{Pipeline} (trait objects) & Yes & --- \\
linfa & Trait-based chaining & --- & --- \\
smartcore & Manual only & --- & --- \\
scikit-learn & \texttt{Pipeline}, \texttt{make\_pipeline} & Yes & Yes \\
\bottomrule
\end{tabular}
\end{center}

scry-learn's \texttt{Pipeline} uses trait objects (\texttt{Box<dyn TransformerBox>} + \texttt{Box<dyn PipelineModel>}) for heterogeneous composition. The \texttt{impl\_pipeline\_model!} macro implements \texttt{PipelineModel} for all 21+ model types.


%=============================================================================
\section{Model Selection and Evaluation}
%=============================================================================

\subsection{Cross-Validation}

{\small
\begin{center}
\begin{tabular}{lcccc}
\toprule
\textbf{Strategy} & \textbf{scry-learn} & \textbf{linfa} & \textbf{smartcore} & \textbf{sklearn} \\
\midrule
train\_test\_split & Yes & Yes & --- & Yes \\
Stratified split & Yes & --- & --- & Yes \\
K-Fold & Yes & Yes & --- & Yes \\
Stratified K-Fold & Yes & --- & --- & Yes \\
Group K-Fold & Yes & --- & --- & Yes \\
Time Series Split & Yes & --- & --- & Yes \\
Repeated K-Fold & Yes & --- & --- & Yes \\
cross\_val\_score & Yes & Yes & --- & Yes \\
cross\_val\_predict & Yes & --- & --- & Yes \\
\bottomrule
\end{tabular}
\end{center}
}

scry-learn implements 7 splitting strategies, closely mirroring sklearn's \texttt{model\_selection} module. All splitters materialize folds eagerly as \texttt{Vec<(Dataset, Dataset)>} (sklearn uses lazy generators yielding index arrays).

%-----------------------------------------------------------------------------
\subsection{Hyperparameter Search}

\begin{center}
\begin{tabular}{lcccc}
\toprule
\textbf{Method} & \textbf{scry-learn} & \textbf{linfa} & \textbf{smartcore} & \textbf{sklearn} \\
\midrule
GridSearchCV & Yes & --- & --- & Yes \\
RandomizedSearchCV & Yes & --- & --- & Yes \\
BayesSearchCV & Yes & --- & --- & Yes (scikit-optimize) \\
HalvingGridSearchCV & --- & --- & --- & Yes \\
\bottomrule
\end{tabular}
\end{center}

scry-learn's hyperparameter search uses the \texttt{Tunable} trait:

\begin{lstlisting}[language=Rust]
pub trait Tunable {
    fn set_param(&mut self, name: &str, value: ParamValue) -> Result<()>;
}
\end{lstlisting}

%-----------------------------------------------------------------------------
\subsection{Metrics}

\begin{center}
\begin{tabular}{lp{3.5cm}p{2cm}p{1.5cm}p{2cm}}
\toprule
\textbf{Category} & \textbf{scry-learn} & \textbf{linfa} & \textbf{smartcore} & \textbf{sklearn} \\
\midrule
Classification & accuracy, balanced\_acc, precision, recall, f1, kappa, log\_loss, confusion matrix, ROC AUC, PR curve & accuracy, confusion matrix & accuracy & 30+ \\
Regression & $R^2$, MSE, MAPE, explained variance & $R^2$, MSE & $R^2$, MSE & 15+ \\
Clustering & ARI, silhouette, CH, DB & --- & --- & 10+ \\
\bottomrule
\end{tabular}
\end{center}


%=============================================================================
\section{Serialization and Versioning}
%=============================================================================

\begin{center}
\begin{tabular}{llll}
\toprule
\textbf{Library} & \textbf{Serialization} & \textbf{Model Versioning} & \textbf{Cross-Language} \\
\midrule
scry-learn & serde (feature-gated) & \texttt{\_schema\_version} field & ONNX export \\
linfa & serde (partial) & --- & --- \\
smartcore & serde (always-on) & --- & --- \\
scikit-learn & pickle/joblib & --- & ONNX (skl2onnx) \\
\bottomrule
\end{tabular}
\end{center}

scry-learn's model versioning prevents stale model deserialization:

\begin{lstlisting}[language=Rust]
// Every model struct has:
#[cfg_attr(feature = "serde", serde(default))]
_schema_version: u32,

// Checked at predict time:
pub fn predict(&self, features: &[Vec<f64>]) -> Result<Vec<f64>> {
    crate::version::check_schema_version(self._schema_version)?;
    // ...
}
\end{lstlisting}

If the serialized model's schema version doesn't match the current library version, prediction returns \texttt{SchemaVersionMismatch} instead of silently producing incorrect results.


%=============================================================================
\section{Unique Features}
%=============================================================================

\subsection{scry-learn}

\begin{itemize}[itemsep=2pt]
\item \textbf{Terminal visualization}: Confusion matrices, ROC curves, learning curves, feature importance plots via \texttt{scry-chart}
\item \textbf{ONNX export}: Manual protobuf serialization (no external deps)
\item \textbf{SHAP explainability}: TreeSHAP + permutation importance
\item \textbf{Prediction checksums}: FNV-1a hash on f64 bits for cross-machine verification
\item \textbf{NaN/Inf rejection}: \texttt{validate\_finite()} at \texttt{fit()} time
\item \textbf{Schema versioning}: Prevents stale model deserialization
\item \textbf{GPU acceleration}: wgpu compute shaders
\item \textbf{Column-major layout}: Optimized for tree and coordinate descent workloads
\item \textbf{Text processing}: CountVectorizer and TfidfVectorizer with sparse output
\item \texttt{\#![deny(unsafe\_code)]}: No unsafe code in the ML crate
\end{itemize}

\subsection{linfa}

\begin{itemize}[itemsep=2pt]
\item \textbf{Type-state pattern}: Compile-time guarantee against unfitted prediction
\item \textbf{Generic float types}: \texttt{F: Float} supports f32 and f64
\item \textbf{Rich trait hierarchy}: \texttt{Fit}, \texttt{Predict}, \texttt{Transformer} for generic composition
\item \textbf{ndarray ecosystem}: Direct interop with Rust numerical computing
\item \textbf{Modular crates}: Independent versioning per algorithm family
\end{itemize}

\subsection{smartcore}

\begin{itemize}[itemsep=2pt]
\item \textbf{Simple API}: Closest to sklearn's feel among Rust libraries
\item \textbf{Multiple matrix backends}: \texttt{BaseMatrix} accepts DenseMatrix, ndarray, or nalgebra
\item \textbf{Low learning curve}: Familiar patterns for sklearn users
\item \textbf{Always-on serde}: All models serializable without feature flags
\end{itemize}

\subsection{scikit-learn}

\begin{itemize}[itemsep=2pt]
\item \textbf{Most comprehensive}: 100+ estimators
\item \textbf{Mature ecosystem}: 15+ years of production hardening
\item \textbf{Cython optimizations}: Critical inner loops compiled to C
\item \textbf{Rich pipeline system}: Pipeline, ColumnTransformer, FeatureUnion
\item \textbf{Community}: Largest ML library user base
\end{itemize}


%=============================================================================
\section{Performance Characteristics}
%=============================================================================

\subsection{Memory Layout Impact}

For \textbf{tree-based algorithms}, scry-learn's column-major layout provides an advantage: finding the best split for a feature requires scanning the entire column, which is contiguous in memory. In row-major layouts, this access has stride \texttt{n\_features}, causing cache misses. sklearn mitigates this by internally converting to Fortran-order.

For \textbf{distance-based algorithms} (KNN, K-Means), row-major is naturally faster since computing distance between two samples accesses contiguous memory. scry-learn compensates with the lazy \texttt{row\_major\_cache}.

\subsection{Abstraction Overhead}

\begin{center}
\begin{tabular}{lll}
\toprule
\textbf{Library} & \textbf{Abstraction} & \textbf{Overhead} \\
\midrule
scry-learn & Direct \texttt{Vec<f64>} indexing & Minimal \\
linfa & ndarray views/slicing & Low \\
smartcore & \texttt{BaseMatrix} trait & Moderate (possible dyn dispatch) \\
scikit-learn & numpy C API & Minimal (Cython bypasses Python) \\
\bottomrule
\end{tabular}
\end{center}

\subsection{Solver Complexity}

\begin{center}
\begin{tabular}{lll}
\toprule
\textbf{Operation} & \textbf{scry-learn} & \textbf{Best Alternative} \\
\midrule
Linear regression & $O(nm^2 + m^3)$ Gauss-Jordan & $O(nm^2)$ Cholesky \\
PCA & $O(m^3 \cdot \text{iters})$ Jacobi & $O(nm^2)$ SVD \\
KNN search & $O(n\log n)$ KD-tree & Same (ball tree for high dims) \\
SVM & $O(n^2)$--$O(n^3)$ SMO & Same across implementations \\
\bottomrule
\end{tabular}
\end{center}

scry-learn's PCA via Jacobi rotations avoids external BLAS/LAPACK dependencies but is asymptotically slower than SVD for large feature counts. For typical ML datasets ($m < 1000$), the difference is negligible.

\subsection{Parallelism Utilization}

scry-learn and sklearn parallelize the most expensive operations (ensemble training, cross-validation). linfa offers optional rayon support. smartcore is entirely single-threaded, which is a significant bottleneck for ensemble methods on multi-core systems.

scry-learn's GPU acceleration via wgpu is unique among Rust ML libraries.


%=============================================================================
\section{Algorithm Coverage Summary}
%=============================================================================

\begin{center}
\begin{tabular}{lcccc}
\toprule
\textbf{Category} & \textbf{scry-learn} & \textbf{linfa} & \textbf{smartcore} & \textbf{sklearn} \\
\midrule
Linear models & 5 & 5 & 5 & 10+ \\
Tree ensembles & 8 (incl. Hist-GBT) & 2 & 2 & 8+ \\
SVM & 4 & 1--2 & 2 & 4 \\
Neural networks & 2 + CNN layers & --- & --- & 2 \\
Naive Bayes & 3 & 1 & 2 & 5 \\
Neighbors & 2 & 2 & 2 & 4+ \\
Clustering & 4 & 4 & 2 & 10+ \\
Anomaly detection & 1 & --- & --- & 5+ \\
Ensemble meta-learners & 2 & --- & --- & 5+ \\
Preprocessing & 10 & 3 & 3 & 20+ \\
Text processing & 2 & --- & --- & 5+ \\
Feature selection & 3 & --- & --- & 10+ \\
Explainability & 2 & --- & --- & 2+ \\
Hyperparameter search & 3 & --- & --- & 4+ \\
\midrule
\textbf{Total models} & \textbf{31+} & \textbf{\textasciitilde15} & \textbf{\textasciitilde12} & \textbf{100+} \\
\bottomrule
\end{tabular}
\end{center}


%=============================================================================
\section{Key Architectural Decisions}
%=============================================================================

\begin{center}
\begin{tabular}{lp{4.5cm}p{5.5cm}}
\toprule
\textbf{Decision} & \textbf{scry-learn Choice} & \textbf{Rationale} \\
\midrule
Memory layout & Column-major default & Optimizes tree splits and coordinate descent \\
Matrix deps & No external (ndarray/nalgebra) & Self-contained, minimal dependency tree \\
Numeric type & \texttt{f64} only & Simplifies APIs, eliminates generic noise \\
Unsafe policy & \texttt{\#![deny(unsafe\_code)]} & Correctness-first philosophy \\
Data bundling & \texttt{Dataset} struct & Prevents shape mismatches, carries metadata \\
Config style & Builder pattern & Ergonomic, matches Rust conventions \\
Fit check & Runtime (\texttt{NotFitted}) & Simpler API than type-state, familiar to sklearn users \\
Visualization & Built-in (scry-chart) & Terminal rendering without external tools \\
Versioning & Schema version field & Prevents stale model deserialization \\
\bottomrule
\end{tabular}
\end{center}

\end{document}
